\chapter{Simulation with aegir}
\label{ch:aegir}

\chaptergoals{7 of 11 in Part~II: the first simulation.}
{\cref{ch:geometry,ch:services,ch:generator}: the geometry, the field maps and a checked HNL sample.}
{write and run a simulation workflow, choose the sensitive volumes, read the validation histograms, and reproduce a run exactly.}

\section{What aegir is}

\repo{aegir} is the SHiP simulation: a set of Phlex plugins that generate or
read particles, transport them through the detector with Geant4, and write
the result. A simulation job is a graph of four kinds of node:

\begin{center}
\begin{tikzpicture}[
  every node/.style={font=\sffamily\small},
  box/.style={draw=shipnavy,fill=codebg,rounded corners=2pt,
    text width=35mm,minimum height=12mm,align=center},
  job/.style={box,draw=shipblue,fill=shipblue!6,text width=55mm},
  flow/.style={-{Stealth[length=2mm]},line width=0.7pt},
  uses/.style={flow,dashed,color=shipblue}]
  \node[box] (src) at (0,0) {source\\\scriptsize\texttt{file\_source}};
  \node[box] (g4) at (5.2,0) {transform\\\scriptsize\texttt{geant4\_module}};
  \node[box] (out) at (10.4,0) {observer\\\scriptsize\texttt{sim\_output\_module}};
  \node[job] (geo) at (2.1,2.4) {job provider\\\scriptsize\texttt{geometry\_geomodel\_provider}};
  \node[job] (fld) at (8.4,2.4) {job provider\\\scriptsize\texttt{field\_map\_provider}};
  \draw[flow] (src) -- node[above,font=\scriptsize]{particles} (g4);
  \draw[flow] (g4) -- node[above,font=\scriptsize]{sim\_result} (out);
  \draw[uses] (geo.south) -- ([xshift=-7mm]g4.north);
  \draw[uses] (fld.south) -- ([xshift=7mm]g4.north);
  \draw[flow] (src.south) -- ++(0,-0.8) -|
    node[pos=0.25,below,font=\scriptsize]{generator record + event header} (out.south);
\end{tikzpicture}
\end{center}

Sources work per event; the geometry and the field are providers at the job
level, built once and shared by all events. The plugins available are:

\begin{center}
\small
\begin{tabularx}{\linewidth}{l l L}
\toprule
Plugin & Kind & Purpose \\
\midrule
\texttt{particle\_gun\_source}       & source   & single particles in a cone \\
\texttt{pythia8\_source}             & source   & Pythia\,8 fixed-target collisions \\
\texttt{fixed\_target\_source}       & source   & Pythia\,8 on tungsten, with the interaction point sampled in the target \\
\texttt{genie\_reader\_source}       & source   & pre-generated GENIE neutrino interactions \\
\texttt{file\_source}                & source   & \textbf{events read from a data-model file (our case)} \\
\texttt{geometry\_builtin\_provider} & provider & a test geometry (tungsten block and silicon planes) \\
\texttt{geometry\_gdml\_provider}    & provider & geometry from a GDML file \\
\texttt{geometry\_geomodel\_provider}& provider & \textbf{the SHiP geometry via the geometry service} \\
\texttt{field\_null\_provider}       & provider & no field \\
\texttt{field\_map\_provider}        & provider & \textbf{field maps via the field service} \\
\texttt{geant4\_module}              & transform& Geant4 transport, multi-threaded \\
\texttt{sim\_output\_module}         & observer & RNTuple output and validation histograms \\
\bottomrule
\end{tabularx}
\end{center}

\section{The HNL workflow}

The workflow for this tutorial is \file{code/workflows/hnl\_sim.jsonnet}. The
workflows shipped with \pkg{aegir} (in its \file{workflows/} directory) build
jobs from the blocks in \file{lib.libsonnet}; this one spells everything out
so that you can read the whole job in one place:

\codefile[linerange={20-107}]{jsonnet}{code/workflows/hnl_sim.jsonnet}

\section{Sensitive volumes}
\label{sec:sensitive}

In default scoring mode, Geant4 creates a \texttt{SimHit} for every
positive-energy-deposit step a particle takes inside a
\emph{sensitive} volume. \pkg{aegir} makes a logical volume sensitive when its
name \emph{contains} one of the fragments in \texttt{sensitive\_volumes}, and
records the \emph{index} of that fragment as \field{detectorId}. Three
consequences:

\begin{enumerate}
  \item The order of the list is the detector numbering. \pkg{Shannon} reads
        \field{detectorId} with the data model's \texttt{detector\_id} enum
        (\cref{ch:datamodel}), so the list must be UBT, SBT, straws,
        calorimeter, timing detector, and nothing else.
  \item A fragment selects everything whose name contains it, including
        passive volumes with similar names. In this workflow the UBT entry
        also catches the air envelopes around the tiles, the SBT entry the
        aluminium walls of the cells, and the calorimeter entry the lead and
        iron absorbers (\cref{ex:sensitive}).
  \item Volumes that match no fragment produce no hits at all. If a detector
        is missing from the list, it is silently invisible.
\end{enumerate}

\begin{note}
The \pkg{aegir} workflows that use the SHiP geometry ship with
\texttt{sensitive\_volumes: ['ScoringPlane']}, which matches nothing in the
real detector: they are set up for muon-flux studies at scoring planes, not
for detector hits. The demonstration list in \file{hnl\_sim.jsonnet} preserves Shannon's
numbering, but does not select active material exclusively.
\end{note}

\begin{organisernote}
To our knowledge there is no reference workflow in \pkg{aegir} that maps the
five detectors in the order \pkg{Shannon} expects; the list in
\file{hnl\_sim.jsonnet} was written for this booklet from the volume names in
\repo{Geometry}. Have the \pkg{aegir}/\pkg{Shannon} developers check it, and
consider adding it to \pkg{aegir}'s \file{lib.libsonnet}.
\end{organisernote}

\begin{pitfall}
The broad UBT, SBT and calorimeter prefixes also record passive-material
energy deposits. Shannon smears them as if they were detector hits; the
example analysis cannot distinguish them using detector ID alone. The veto
and calorimeter response are therefore teaching proxies, not physical
predictions. A production workflow needs active-volume selection with an
explicit detector-ID mapping, or a geometry-aware hit filter.
\end{pitfall}

\section{Running it}

\begin{exercise}{Simulate 100 HNL decays}
\label{ex:sim100}
Open a shell in the \pkg{aegir} environment and go to your working
directory, which must contain \file{ship\_geometry.db}, the two field maps and
\file{hnl\_events.root}:
\begin{shell}
cd "$SHIP_TUTORIAL/aegir"
pixi shell
cd "$SHIP_TUTORIAL/work"
time phlex -j 4 -c <(jsonnet \
    --ext-str events=100 \
    --ext-str infile=hnl_events.root \
    --ext-str simout=hnl_sim.root \
    --ext-str histo=hnl_sim_validation.root \
    $CODE/workflows/hnl_sim.jsonnet)
\end{shell}
Watch the log: the geometry is read and converted, the field maps are loaded
and attached to their volumes, the Geant4 seed is printed, and then the event
loop reports progress every 100 events.
\end{exercise}

\begin{checkpoint}
Two new files, \file{hnl\_sim.root} and \file{hnl\_sim\_validation.root}, and
no error. With the Python lines from \cref{ex:openfile},
\file{hnl\_sim.root} shows 100 entries in the \texttt{events} RNTuple with the
five fields of \cref{sec:filelayout}.
\end{checkpoint}

\begin{pitfall}
The driver creates exactly \texttt{events} events whether or not the input has
that many entries. Past the end of the file, \texttt{file\_source} publishes
\emph{empty} events with weight~1. Always ask for at most as many events as
the input holds (\cref{ex:checkinput} printed the number).
\end{pitfall}

\begin{exercise}{Read the validation histograms}
\pkg{aegir} fills a few histograms as it runs. Open them with a ROOT browser:
\begin{shell}
root -l hnl_sim_validation.root
root [1] new TBrowser
\end{shell}
\begin{enumerate}
  \item \texttt{h\_hit\_detector}: which detector IDs occur, and roughly in
        which proportions? Is any of the five missing?
  \item \texttt{h\_hit\_z}: identify the straw stations, the timing detector
        and the calorimeter. Are there hits upstream of the trackers, and from
        which detector?
  \item \texttt{h\_mc\_multiplicity}: every input event holds three
        generator particles. Why does this histogram say two?
\end{enumerate}
\end{exercise}

\begin{solution}
\begin{enumerate}
  \item IDs 2 (straws), 3 (calorimeter) and 4 (timing) dominate; 1 (SBT)
        appears when a daughter or a secondary leaves through the vessel
        walls. ID 0 (UBT) should be essentially absent: the UBT is upstream of
        the decay volume.
  \item Four clusters between 84 and 95.1\,m for the straws, one at about
        96\,m for the timing detector, 96.9--99.8\,m for the calorimeter.
        Hits upstream come from the SBT along the decay volume.
  \item The validation histograms count final-state (status~1) entries only,
        the ones Geant4 tracks; the decayed HNL is not one of them.
\end{enumerate}
\end{solution}


\section{Reproducibility and threads}

\begin{exercise}{Same seed, same events}
Rename your output, run the same command again, and compare selected per-event summaries
with the summary script from \cref{ch:analysis}, or simply compare the
validation histograms. Then run once more with \cmd{-j 1} and once with
\cmd{--tla-code seed=1} added to the \cmd{jsonnet} command.
\begin{shell}
python $CODE/analysis/sim_summary.py hnl_sim.root hnl_sim_run2.root
\end{shell}
Which summaries agree exactly, and which differ? Equal summaries are a
smoke check, not proof that every field in the two files is identical.
\end{exercise}

\begin{solution}
For a fixed software build and identical inputs, the expected result is
matching event content with the default seed, also with \cmd{-j 1}:
Geant4 is reseeded per event from (seed, event index), so neither the number
of threads nor the order in which events finish changes any event. The
entries can still be in a different order in the file, which is why
\file{sim\_summary.py} compares per-event quantities keyed by
\field{event\_number}. The run with \cmd{seed=1} differs event by event but
should be statistically compatible; a small sample is not a precision test.
\end{solution}

\section{Going further}

\begin{exercise}[colback=white]{Optional: switch the field off}
Make a copy of the workflow that uses \texttt{field\_null\_provider} (no
configuration keys) instead of \texttt{field\_map\_provider}, simulate the
same events into \file{hnl\_sim\_nofield.root}, and keep the file: you will
look at the difference in the event display (\cref{ch:display}) and in the
analysis (\cref{ch:analysis}).
\end{exercise}

\begin{exercise}[colback=white]{Optional: full statistics}
Simulate the whole input file in the background while you continue with the
next chapters. With \cmd{-j} set to the number of cores of your machine this
takes \tbd{time for the provided sample}.
\end{exercise}
